\section{Inference-Time Reasoning：不改权重也能改变计算图}

有了 pretrained model，下一步不一定立刻更新参数。Prompt、search、decomposition、code execution 和 note insertion 都能在 inference time 改变问题求解过程。课程把这些方法放在同一条轴上：它们给模型更多中间计算或外部结构，但成本、可验证性和错误传播模式不同。


\lecturefigure{slide-063.jpg}{Post-training 或 inference-time compute：不同干预位置}{CS25 V5 official deck, p.~63}

\teachervoice{00:25:09--00:26:20，Chelsea 将 fine-tuning、feedback、prompting 与 retrieval 并列，并指出 CoT 提供一扇可观察窗口。讲义提醒：可读 rationale 不必等于模型真正的因果计算过程。}

\subsection{Linear Chain-of-Thought 与示例}

上一小节把 post-training 分成多种接口；这里先研究最轻量的一类：不改参数，只要求模型输出中间步骤。读 CoT 例子时应分别检查 problem decomposition、每一步局部有效性和 final answer，不能因为文字更长就认定推理更深。还要记录 token budget，因为 direct answer 与长 chain 的计算量并不相同。

CoT 把直接映射 $x\rightarrow y$ 改为 $x\rightarrow z_{1:k}\rightarrow y$，其中 $z_i$ 是中间文本步骤。其收益可能来自把复杂函数分解为短局部变换、增加 test-time compute，或使错误更容易定位；但中间步骤也可能流畅而错误。


\lecturefigure{slide-064.jpg}{Chain-of-Thought：把复杂推理展开为可见中间步骤}{CS25 V5 official deck, p.~64}


\lecturefigure{slide-065.jpg}{CoT 示例：直接回答与逐步推理的对照}{CS25 V5 official deck, p.~65}

\begin{warningbox}{Faithfulness 与 usefulness 要分开}
CoT 可以提高正确率或可调试性，却不保证每句话忠实反映内部机制。评估至少要同时看 final accuracy、step validity、self-consistency 与对无关 rationale 的敏感度。
\end{warningbox}

\subsection{怎样改进 chain：搜索、监督与小模型}

改进方法包括让模型回顾前文、搜索多条路径、用 verifier 评分、训练 step-level supervision，或从更强模型蒸馏 reasoning traces。对小模型，长 chain 可能超出容量，因此 compact decomposition 与 semantic understanding 更关键。


\lecturefigure{slide-066.jpg}{Improving CoT：verifier、self-consistency 与过程监督}{CS25 V5 official deck, p.~66}


\lecturefigure{slide-067.jpg}{Smaller models：逐步提示、语义理解与容量限制}{CS25 V5 official deck, p.~67}


\lecturefigure{slide-068.jpg}{Generalizing CoT：跨领域、跨难度与结构迁移}{CS25 V5 official deck, p.~68}

\begin{knowledgebox}{深读：CoT generalization 至少有四种含义}
一是同分布新题，模型复用相同解题模板；二是难度外推，题目需要更多步骤；三是领域迁移，从数学 chain 迁到常识、规划或代码；四是结构迁移，从示例中的表面格式抽象出可组合策略。Slides 66--68 的方法可能改善其中一类而伤害另一类。评估时应改变措辞、顺序和无关信息，测试模型是否依赖 template；用 adversarial intermediate step 检查它是否真正读取 chain；比较 small/large models 判断收益是否受容量限制；最后在没有示例的任务上测 transfer。所谓“generalized CoT”只有在这些边界被说明后才有清楚含义。
\end{knowledgebox}

\teachervoice{00:26:25--00:29:12，课堂把 CoT 扩展为一组 inference procedures，而不是单一 prompt。实践经验：比较方法时要固定 token budget 或 wall-clock，否则“更会推理”可能只是“用了更多采样”。}

\subsection{Tree、Program 与 decomposition}

Tree of Thoughts 在每一步生成多个候选并搜索：
\[
s_{t+1}^{(j)}=f_{\theta}(s_t,a_t^{(j)}),\qquad
j^{\star}=\arg\max_j V(s_{t+1}^{(j)}).
\]
$s_t$ 是当前状态，$a_t^{(j)}$ 是候选 thought，$V$ 是评估器。它把生成变成 search，但性能取决于 proposal diversity 和 evaluator calibration。


\lecturefigure{slide-069.jpg}{Tree-of-Thoughts：多分支生成、评估与回溯}{CS25 V5 official deck, p.~69}


\lecturefigure{slide-070.jpg}{Program-of-Thought：让代码解释器承担精确计算}{CS25 V5 official deck, p.~70}

Socratic decomposition、least-to-most 与 computation graph 都把问题拆成依赖子问题；区别在于拆解是对话式、顺序式还是显式图结构。Program-of-Thought 把算术交给 interpreter，降低语言模型在精确执行上的负担，却引入 sandbox 与代码错误处理需求。


\lecturefigure{slide-071.jpg}{Socratic questioning：递归追问关键子问题}{CS25 V5 official deck, p.~71}


\lecturefigure{slide-072.jpg}{Least-to-Most prompting：先解简单子问题再组合}{CS25 V5 official deck, p.~72}


\lecturefigure{slide-073.jpg}{Program of Thoughts：语言规划与可执行计算分工}{CS25 V5 official deck, p.~73}


\lecturefigure{slide-074.jpg}{Computation Graphs：把依赖关系显式化为图}{CS25 V5 official deck, p.~74}


\lecturefigure{slide-075.jpg}{Self-Notes：在输入序列中插入可检索的中间注释}{CS25 V5 official deck, p.~75}

\begin{knowledgebox}{深读：Reasoning method 的验收不能只看 final accuracy}
对线性 CoT，应抽样检查每一步是否合法以及错误是否被后续偶然抵消；对 tree search，要测 proposal diversity、branching factor、evaluator accuracy 和搜索成本；对 Program-of-Thought，要把 planning error 与 execution error 分开，并在 sandbox 中限制文件、网络和运行时间；对 decomposition，要检查子问题是否覆盖原目标、依赖顺序是否完整；对 self-notes，要观察插入内容是否真的被后续 token 使用。还应比较相同总 token/tool budget 下的 baseline，因为多采样本身就能提高 pass rate。最终报告至少包含 accuracy、cost、latency、verification rate、abstention 与 failure taxonomy。这样才能判断方法增加的是可验证计算，还是只是更长、更昂贵的生成。
\end{knowledgebox}


\lecturefigure{slide-076.jpg}{Latent-space CoT：在隐藏空间增加中间计算}{CS25 V5 official deck, p.~76}

\begin{knowledgebox}{读图：reasoning taxonomy 的两个坐标轴}
横向可按\textbf{表示空间}分为文本 chain、程序、图和 latent state；纵向可按\textbf{控制策略}分为单路径生成、多样本投票、树搜索与外部 verifier。Slide 76 把思考放入 latent space，减少文字带宽，却牺牲可检查性。前面 slides 则用更显式结构换取验证能力。比较方法时应固定 total sampled tokens、tool calls 和 evaluator compute，否则搜索方法的收益无法与基础模型能力分离。
\end{knowledgebox}

\begin{table}[H]
\centering
\small
\begin{tabularx}{\textwidth}{L{0.18\textwidth} X X X}
\toprule
方法 & 增加的资源 & 优势 & 主要失败点 \\
\midrule
Linear CoT & 额外生成 token & 简单、可读 & 错误链条、自信幻觉 \\
Tree search & 多样本与 evaluator & 可回溯、多路径 & 成本高、评分器偏差 \\
Program/Tool & 外部执行器 & 精确计算、可验证 & 安全、接口和代码错误 \\
Decomposition & 子问题结构 & 降低局部难度 & 错误拆分、依赖遗漏 \\
Latent reasoning & 隐状态计算 & 不受文字带宽限制 & 难解释、难监督 \\
\bottomrule
\end{tabularx}
\caption{Inference-time reasoning 方法按资源与失败模式分类。}
\end{table}

\begin{lstlisting}[caption={预算受限的 reasoning search 伪代码}]
frontier = [initial_state(problem)]
for depth in range(max_depth):
    candidates = expand(frontier, samples_per_state)
    checked = verify_with_tools(candidates)
    frontier = top_k(checked, score="validity_then_reward")
    if solved(frontier[0]):
        return frontier[0]
return best_with_uncertainty(frontier)
\end{lstlisting}

\subsection{本章小结}

Inference-time reasoning 改变的是求解计算图，不必更新 pretrained weights。CoT、tree、program、decomposition 与 latent chain 分别增加 token、search、工具、结构或隐状态计算；公平比较必须控制预算并报告 verifier、工具和失败恢复。

